% Generated draft: requires mathematical and source review.
\documentclass[11pt]{article}
\usepackage[margin=25mm]{geometry}
\usepackage{fontspec}
\setmainfont{Latin Modern Roman}
\usepackage{amsmath,amssymb,mathtools,unicode-math}
\setmathfont{Latin Modern Math}
\usepackage{enumitem,needspace,xurl,hyperref}
\setlength{\parindent}{0pt}
\setlength{\parskip}{5pt}
\setlength{\emergencystretch}{3em}
\newcommand{\N}{\mathbb{N}}
\newcommand{\Z}{\mathbb{Z}}
\newcommand{\Q}{\mathbb{Q}}
\newcommand{\R}{\mathbb{R}}
\newcommand{\C}{\mathbb{C}}
\newcommand{\F}{\mathbb{F}}
\newcommand{\A}{\mathbb{A}}
\newcommand{\PP}{\mathbb{P}}
\newcommand{\E}{\mathbb{E}}
\newcommand{\eps}{\varepsilon}
\newcommand{\dd}{\,\mathrm{d}}
\begin{document}
\section*{UnsolvedMath ID: 20001991}
\textbf{Source problem code:} catalogue entry
\subsection{Definitions and mathematical statement}
\textbf{Question A.9, Ilia Itenberg.} What restrictions and bounds can be given for the Betti numbers of real tropical varieties? The source relates real tropical hypersurfaces to $T$-hypersurfaces from combinatorial patchworking and asks to reinterpret known topological restrictions tropically. It specifies neither a class of varieties nor compactness, smoothness, coefficient field or a proposed universal bound; this is a research question rather than a single numerical conjecture.

\textbf{Supplied nonsingular compact patchwork setting.} Let $X$ be a compact nonsingular tropical variety of dimension $d$ with real phase/patchworking data producing a compact real patchwork $X_{\R}$. Assume the relevant Renaudineau--Shaw or Rau--Renaudineau--Shaw filtered-chain spectral sequence exists and converges to $H_*(X_{\R};\F_2)$. Let $\mathcal F_p$ be its tropical multi-tangent coefficient systems and put
\[
 U_q=\sum_p\dim_{\F_2}H_q(X;\mathcal F_p),\qquad
 b_q=\dim_{\F_2}H_q(X_{\R};\F_2),\qquad D_q=U_q-b_q.
\]
The precise realization hypotheses and tropical coefficient conventions must be checked against the cited spectral sequence. In particular these bounds are not asserted for every singular or noncompact real tropical variety.

Index each spectral-sequence page by filtration and \emph{total homological degree} $q$, so every differential lowers $q$ by one; $q$ is not silently treated as the second coordinate of a conventional total-degree sum. Let $a_q$ be the sum, over all pages and filtration degrees, of the ranks of differentials with source in degree $q$. For this finite-dimensional bounded spectral sequence, the aid asserts
\[
 D_q=a_q+a_{q+1},\qquad a_q\in\Z_{\ge0},\qquad
 a_0=a_{d+1}=0,\quad 0\le q\le d.
\]
Consequently $D_q\ge0$, and more precisely
\[
 \sum_{j=0}^q(-1)^{q-j}D_j=a_{q+1}\ge0
 \quad(0\le q<d),\qquad
 \sum_{j=0}^d(-1)^jD_j=0,
 \qquad \sum_{q=0}^dD_q=2\sum_{q=1}^d a_q.
\]
Thus the degreewise deficits have nonnegative alternating prefixes and the total deficit is exactly twice the total rank of spectral cancellations. These are necessary constraints on an actual patchwork; the aid does not claim that every vector satisfying them is geometrically realizable. The individual bounds $b_q\le U_q$ and Euler equality alone do not encode all the prefix tests in general dimension.

For curves, $D_0=D_1$; in the usual connected compact smooth genus-$g$ situation with $U_0=g+1$, this is the Harnack component defect $g+1-b_0$. For compact closed surfaces with the required mod-two Poincar\'e symmetries $U_2=U_0$ and $b_2=b_0$, one obtains $D_2=D_0$ and $D_1=2D_0$. These low-dimensional statements retain their compactness and duality qualifications. The aid supplies restrictions from a spectral sequence, not a classification of real tropical Betti vectors or an answer for the full scope of Question A.9.
\subsection{Short English statement}
Itenberg asks broadly about Betti numbers of real tropical varieties. For the compact nonsingular patchworks covered by the supplied spectral sequence, the aid expresses each Betti deficit as the sum of outgoing and incoming differential ranks. This yields alternating-prefix restrictions and low-dimensional defect identities, without classifying all possible Betti numbers.
\subsection{Sources}
\begin{itemize}
\item Local UnsolvedMath registry, permanent ID 20001991, source code catalogue entry; canonical \texttt{aim-geometry-notes.json}, notes[328].
\item AIM \emph{Amoebas and tropical geometry}, item A.9, contributed by Ilia Itenberg: \url{https://aimath.org/WWN/amoebas/amoebas.pdf}. The record reports agreement between its HTML and PDF sources; no new verification is claimed here.
\item The research aid names the Renaudineau--Shaw and Rau--Renaudineau--Shaw spectral sequences, but supplies no full titles, dates, theorem labels or URLs. No additional citations are invented.
\item Registry research attempt 1, dated 2026-08-14: partial result, rigor ``proved'', novelty confidence low. The machine-generated, non-peer-reviewed aid requires independent verification before citation.
\end{itemize}
\end{document}
